\documentclass[12pt]{article}
\usepackage[utf8]{inputenc}
\usepackage[T1]{fontenc}
\usepackage{hyperref}
\title{Ngarinyin Language Resources for K-12 Curriculum Development}
\author{Compiled by Claude Sonnet 5 (claude-sonnet-5) for Andr\'as Kornai}
\date{}
\begin{document}
\maketitle
\section{Introduction}
Ngarinyin (ISO 639-3 ung; Glottocode ngar1284) is a Worrorran language of the Kimberley region, Western Australia. It is severely endangered, with very few elderly speakers. My search found no school curriculum materials and no published dictionary.
\section{Primary Educational Resources}
None found.
\section{Linguistic and Reference Materials}
Overview: \url{https://en.wikipedia.org/wiki/Ngarinyin_language}. The AIATSIS collection page for Ngarinyin (K18), which lists earlier grammar work such as Capell (1940), Coate \& Oates (1970) and Rumsey (1978, 1982): \url{https://collection.aiatsis.gov.au/node/1508631}. The Endangered Languages Project entry: \url{https://endangeredlanguages.com/node/37772}.
\section{Religious and Community Resources}
None verified.
\section{Digital and Audio Resources}
The ELDP-funded project ``Documentation and description of Ngarinyin'' (Stef Spronck, ANU, 2011--2012) deposited elicited and spontaneous speech recordings: \url{https://eldp.access.preservica.com/dk0121/}. Project record: \url{https://researchportalplus.anu.edu.au/en/projects/documentation-and-description-of-ngarinyin-a-language-of-the-kimb/}.
\section{Curriculum Development Context}
The language is not known to be taught in schools. Any curriculum would have to be built from archival recordings and older grammars, in consultation with the Ngarinyin community. No contacts are listed here.
\end{document}